% Complete unabridged proof from architecture/proofs/ONE_ROW_SEMANTIC_ANCHOR_PROOF.md
\section{One semantic Pell anchor and an aligned amplification index}



\subsection{Result and meaning of the equation}



The main new natural-domain construction is the sum of squares of the 179 quadratic integer residuals in \texttt{strict\_\allowbreak{}natural229/\allowbreak{}natural/\allowbreak{}equation.json}. It has 229 nonnegative integer variables. Its independently collected quartic has 9,184 nonzero monomials and coefficient height 83,954,634. The sparse companion \texttt{strict\_\allowbreak{}natural230} has 230 variables, 180 rows, 6,858 monomials and height 83,954,472. An additional exact-projection companion has 228 variables, 178 rows, 9,866 monomials and height 2,445,412,949,744. It trades much greater coefficient height for one fewer variable and is not the recommended balanced form. Each coefficient is supplied literally in the respective \texttt{expanded\_\allowbreak{}polynomial.json}; no exponentiation, factorial, binomial, bounded quantifier, or implicit sequence is a primitive of that polynomial.



The existence claim is unchanged: with the inherited ordinary numerical C222 calculus calibration, the polynomial has a natural zero exactly when ZFC has a finite syntactic refutation. Consequently ZFC proves that its no-zero sentence is equivalent to Con(ZFC). This is neither a proof of consistency, a global minimality assertion, nor a newly checked end-to-end proof-assistant or object-language ZFC derivation.



There are two new compatible reductions. First, a single bounded-error row reuses the already required derived-label Pell index as the full anchor. Its index need not be even and the size frame need not be doubled. Second, the amplification exponent is the largest required quotient-range index x, rather than a separate large affine expression. The two beta quotient ranges are sharpened from u,v\textless{}\allowbreak{}=\allowbreak{}x to u,v\textless{}\allowbreak{}x, as their defining equalities already permit. The incoming factorial/\allowbreak{}product fusion, modular binomial graphs, dominant denominator, consecutive Pell coordinates, and self-sized radix are explicitly retained and credited.



\subsection{A general one-row semantic-anchor lemma}



Write X\_\allowbreak{}j(A),Y\_\allowbreak{}j(A) for the ordinary Pell sequences defined by



\begin{verbatim}
(A+sqrt(A^2-1))^j = X_j(A)+Y_j(A)sqrt(A^2-1).
\end{verbatim}

Assume integers W\textgreater{}\allowbreak{}=\allowbreak{}3, 2\textless{}\allowbreak{}=\allowbreak{}H\textless{}\allowbreak{}W, A-1\textgreater{}\allowbreak{}W\textasciicircum{}\allowbreak{}3, r\textless{}\allowbreak{}=\allowbreak{}W-2 and requested indices 1\textless{}\allowbreak{}=\allowbreak{}k\_\allowbreak{}i\textless{}\allowbreak{}W. Retain a complete natural indexed-Pell certificate at index H, with principal pair (U,V), omitting only its separate inequality H\textless{}\allowbreak{}=\allowbreak{}V. For each i retain a natural principal Pell norm. Impose the single row



\begin{verbatim}
V = H + sum_i (Y_i-k_i)^2 + (A-1)Q,       Q>=0.             (1)
\end{verbatim}

Then the principal pair is (X\_\allowbreak{}(k\_\allowbreak{}i)(A),Y\_\allowbreak{}(k\_\allowbreak{}i)(A)) for every i. Conversely, canonical pairs extend to (1) whenever H\textgreater{}\allowbreak{}=\allowbreak{}2 max\_\allowbreak{}i k\_\allowbreak{}i. All private witnesses of the complete indexed-Pell certificate retain their inherited natural domains. The assertion does not license making any private witness signed.



For soundness, (1) supplies V\textgreater{}\allowbreak{}=\allowbreak{}H, so the complete anchor theorem first gives V=\allowbreak{}Y\_\allowbreak{}H(A). Since H\textgreater{}\allowbreak{}=\allowbreak{}2, V\textgreater{}\allowbreak{}=\allowbreak{}Y\_\allowbreak{}2(A)=\allowbreak{}2A\textgreater{}\allowbreak{}2W. Classification of each principal norm gives Y\_\allowbreak{}i=\allowbreak{}Y\_\allowbreak{}(j\_\allowbreak{}i)(A), j\_\allowbreak{}i\textgreater{}\allowbreak{}=\allowbreak{}0. If j\_\allowbreak{}i\textgreater{}\allowbreak{}=\allowbreak{}H, then Y\_\allowbreak{}i\textgreater{}\allowbreak{}=\allowbreak{}V and



\begin{verbatim}
(Y_i-k_i)^2 >= (V-W+1)^2 > V.
\end{verbatim}

For the last strict inequality, V\textgreater{}\allowbreak{}=\allowbreak{}2W gives V-W+\allowbreak{}1\textgreater{}\allowbreak{}=\allowbreak{}V/\allowbreak{}2+\allowbreak{}1, whose square is larger than V. This contradicts (1), which bounds the sum of error squares by V-H\textless{}\allowbreak{}V. Hence j\_\allowbreak{}i\textless{}\allowbreak{}H\textless{}\allowbreak{}W. The Pell recurrence modulo A-1 gives Y\_\allowbreak{}j(A)=\allowbreak{}j modulo A-1. Reducing (1) therefore gives



\begin{verbatim}
A-1 divides E = sum_i (j_i-k_i)^2,
0<=E<=r(W-1)^2<W^3<A-1.
\end{verbatim}

Thus E=\allowbreak{}0 and j\_\allowbreak{}i=\allowbreak{}k\_\allowbreak{}i. In particular, no separate hypothesis k\_\allowbreak{}i\textless{}\allowbreak{}H was used in soundness. If ordinary coordinates are used, natural X\_\allowbreak{}i supplies the positive abscissa. Consecutive coordinates are handled below.



For completeness, put s=\allowbreak{}max\_\allowbreak{}i k\_\allowbreak{}i\textgreater{}\allowbreak{}=\allowbreak{}1 and z=\allowbreak{}Y\_\allowbreak{}s(A)\textgreater{}\allowbreak{}=\allowbreak{}1. If H\textgreater{}\allowbreak{}=\allowbreak{}2s, then



\begin{verbatim}
V=Y_H(A)>=Y_(2s)(A)=2X_s(A)Y_s(A)>2(A-1)z^2.
\end{verbatim}

As Y\_\allowbreak{}(k\_\allowbreak{}i)(A)\textgreater{}\allowbreak{}=\allowbreak{}k\_\allowbreak{}i, each actual square error is at most z\textasciicircum{}\allowbreak{}2. Therefore



\begin{verbatim}
V-sum_i(Y_(k_i)(A)-k_i)^2 > [2(A-1)-r]z^2 > H,
\end{verbatim}

where A-1\textgreater{}\allowbreak{}W\textasciicircum{}\allowbreak{}3, r\textless{}\allowbreak{}=\allowbreak{}W-2 and H\textless{}\allowbreak{}W imply 2(A-1)-r\textgreater{}\allowbreak{}H. The numerator V-H-sum\_\allowbreak{}i(Y\_\allowbreak{}(k\_\allowbreak{}i)(A)-k\_\allowbreak{}i)\textasciicircum{}\allowbreak{}2 is positive and divisible by A-1. Its quotient is the required natural Q. The inherited indexed-Pell completeness theorem supplies fresh private anchor witnesses at this A and H. This proof uses neither parity of H nor a doubled frame.



\subsection{The literal semantic frame and anchor certificate}



Let t=\allowbreak{}t0+\allowbreak{}12 and take the unchanged positive encoded representatives C,A\_\allowbreak{}a,A\_\allowbreak{}b,A\_\allowbreak{}c,A\_\allowbreak{}d,A\_\allowbreak{}S,A\_\allowbreak{}Sp,A\_\allowbreak{}u,A\_\allowbreak{}v. The derived nonnegative label polynomial J is the sum of the nine axiom conclusion polynomials and a,b,I(a,b),A(a,b),1, with the same Cantor constructors and arithmetic source as the predecessor. In particular J\textgreater{}\allowbreak{}=\allowbreak{}I(A\_\allowbreak{}a,A\_\allowbreak{}b)\textgreater{}\allowbreak{}=\allowbreak{}2A\_\allowbreak{}a\textasciicircum{}\allowbreak{}2\textgreater{}\allowbreak{}=\allowbreak{}2. Use H=\allowbreak{}J in (1).



The frame is



\begin{verbatim}
W = 15 + sum(distinct exact-power bases, exponents and outputs
             together with all modular numerator indices).
\end{verbatim}

There are thirteen independent seeds; J itself uses the full anchor pair and is not also a seed. Thus H\textless{}\allowbreak{}W, every intended index is \textless{}\allowbreak{}W and r=\allowbreak{}13\textless{}\allowbreak{}=\allowbreak{}W-2. The current code computes the offset as the actual seed count plus two, rather than trusting an unattested hard-coded census.



The scale rows are



\begin{verbatim}
L^2-((W+1)^2-1)(WZ)^2=1,    Z=z+1,    A=L+sigma,
\end{verbatim}

with all coordinates natural. Pell classification gives L=\allowbreak{}X\_\allowbreak{}e(W+\allowbreak{}1), WZ=\allowbreak{}Y\_\allowbreak{}e(W+\allowbreak{}1) with e\textgreater{}\allowbreak{}0. Reduction modulo W gives W|e, so e\textgreater{}\allowbreak{}=\allowbreak{}W and A\textgreater{}\allowbreak{}=\allowbreak{}L\textgreater{}\allowbreak{}=\allowbreak{}X\_\allowbreak{}W(W+\allowbreak{}1)\textgreater{}\allowbreak{}=\allowbreak{}(W+\allowbreak{}1)\textasciicircum{}\allowbreak{}W. In particular A-1\textgreater{}\allowbreak{}W\textasciicircum{}\allowbreak{}3 and A\textgreater{}\allowbreak{}W+\allowbreak{}1, before interpreting any exact power. The free sigma permits choosing A in the later common-root congruence class.



The full anchor is exactly the inherited three-norm indexed-Pell graph. Using private names U1,V1,B1,S1,R1,d,e,f,g,h, it is



\begin{verbatim}
U^2-(A^2-1)V^2=1,
U1^2-(A^2-1)V1^2=1,
S1^2-(B1^2-1)R1^2=1,
B1=1+4V(d+1),                 V1=V^2(e+1),
B1=A+U1*f,                   S1=U+U1*g,
R1=H+4V*h.
\end{verbatim}

Its omitted H\textless{}\allowbreak{}=\allowbreak{}V condition follows from (1). The existing constructive criterion also supplies the nonnegative oriented congruence quotients: H\textgreater{}\allowbreak{}=\allowbreak{}2 gives V\textgreater{}\allowbreak{}=\allowbreak{}2A; B1\textgreater{}\allowbreak{}1 and B1=\allowbreak{}1 modulo 4V give B1\textgreater{}\allowbreak{}A; R1\textgreater{}\allowbreak{}=\allowbreak{}H\textgreater{}\allowbreak{}=\allowbreak{}2 and Pell classification at B1 give R1\textgreater{}\allowbreak{}=\allowbreak{}2B1, hence S1\textgreater{}\allowbreak{}U. No unsupported new indexed-Pell criterion is being substituted.



\subsection{Power recovery and modular numerators}



Each exact request m=\allowbreak{}n\textasciicircum{}\allowbreak{}k retains



\begin{verbatim}
T_n=2An-n^2-1,
X_k(A)+nY_k(A)=A Y_k(A)+m+T_n*q_n,
\end{verbatim}

with T\_\allowbreak{}n,q\_\allowbreak{}n natural. The first row excludes n=\allowbreak{}0. Every requested exponent is positive from the raw expressions, endpoints and positive representative offsets. After (1), the ordinary Pell congruence gives m=\allowbreak{}n\textasciicircum{}\allowbreak{}k modulo T\_\allowbreak{}n. The frame bounds n,k,m by W, while A\textgreater{}\allowbreak{}=\allowbreak{}(W+\allowbreak{}1)\textasciicircum{}\allowbreak{}W gives T\_\allowbreak{}n\textgreater{}\allowbreak{}W\textasciicircum{}\allowbreak{}W+\allowbreak{}W. Thus the congruence is an equality. Canonical positive powers provide nonnegative quotients by the inherited positive-power argument.



The six exact requests are



\begin{verbatim}
V0=(F^2+1)^t,       N=B^K,           q=(N*S)^K,
D_t=q^t,           D_y=q^(y+1),     D_x=q^x,
K=x,             S=1+C+sum(eight encoded columns).
\end{verbatim}

The extra denominator D\_\allowbreak{}(t+\allowbreak{}1)=\allowbreak{}q\textasteriskcentered{}\allowbreak{}D\_\allowbreak{}t is multiplication, not a seventh exact power. Aligning K with x makes two formerly distinct index requests literally the same atom. All eleven modular numerator indices remain in W.



The common-root equation is



\begin{verbatim}
2A(q+1)=(q+1)^2+1+M*h0,       M=N*D_x.                    (2)
\end{verbatim}

Once the six powers and endpoints are decoded, x\textgreater{}\allowbreak{}t+\allowbreak{}1 and x\textgreater{}\allowbreak{}y+\allowbreak{}1; hence M is divisible by qD\_\allowbreak{}t for the one exact binomial, and by N times every range denominator. Substituting s=\allowbreak{}q+\allowbreak{}1-A into the Pell identity modulo M gives



\begin{verbatim}
X_n(A)+(q+1-A)Y_n(A) = (q+1)^n modulo M.                 (3)
\end{verbatim}

The left side is positive: X\_\allowbreak{}n(A)\textgreater{}\allowbreak{}(A-1)Y\_\allowbreak{}n(A) for n\textgreater{}\allowbreak{}0, and every n here is positive. The literal modular digit equations are the same as the incoming common-root compiler, with Euclidean remainder r\textless{}\allowbreak{}q\textasciicircum{}\allowbreak{}k and, for the exact binomial only, output B\textless{}\allowbreak{}q. Thus they recover exactly the desired low digit or binomial divisibility once the carry bounds below hold.



For completeness, N|q, so all prime divisors of M divide q and gcd(q+\allowbreak{}1,M)=\allowbreak{}1. Hence gcd(2(q+\allowbreak{}1),M) is one or two. In the latter case q+\allowbreak{}1 is odd and (q+\allowbreak{}1)\textasciicircum{}\allowbreak{}2+\allowbreak{}1 is even, so the linear congruence (2) is soluble. Choose A in its solution class arbitrarily large after all semantic data are fixed; both sigma and h0 can be natural. Canonical Pell residues are positive and congruent to their least nonnegative required residues, so all modular digit quotients are natural.



\subsection{Bootstrap and factorial/\allowbreak{}product fusion}



The factorial/\allowbreak{}product output B is shared by the two ceiling rows



\begin{verbatim}
V0+B=B*F+u+1,       B=u+v+1.
\end{verbatim}

They imply B\textgreater{}\allowbreak{}=\allowbreak{}1 and V0+\allowbreak{}v=\allowbreak{}BF. Since V0=\allowbreak{}(F\textasciicircum{}\allowbreak{}2+\allowbreak{}1)\textasciicircum{}\allowbreak{}t\textgreater{}\allowbreak{}=\allowbreak{}1, F\textgreater{}\allowbreak{}=\allowbreak{}1. Also V0\textgreater{}\allowbreak{}F\textasciicircum{}\allowbreak{}2, so B\textgreater{}\allowbreak{}F and V0\textless{}\allowbreak{}=\allowbreak{}BF\textless{}\allowbreak{}B\textasciicircum{}\allowbreak{}2. The raw endpoint x\textgreater{}\allowbreak{}=\allowbreak{}F+\allowbreak{}1 gives K=\allowbreak{}x\textgreater{}\allowbreak{}=\allowbreak{}2. Consequently



\begin{verbatim}
V0<B^2<=N<=q.
\end{verbatim}

For r\textgreater{}\allowbreak{}=\allowbreak{}1 and n\textgreater{}\allowbreak{}=\allowbreak{}1, every coefficient binom(r+\allowbreak{}n,j), 0\textless{}\allowbreak{}=\allowbreak{}j\textless{}\allowbreak{}=\allowbreak{}n, is at most (r+\allowbreak{}1)\textasciicircum{}\allowbreak{}n. This follows by Pascal induction for j\textless{}\allowbreak{}n and the product formula for j=\allowbreak{}n. Apply this with r=\allowbreak{}F\textasciicircum{}\allowbreak{}2,n=\allowbreak{}t. All coefficients through t are \textless{}\allowbreak{}=\allowbreak{}V0\textless{}\allowbreak{}q, so the exact modular digit really gives B=\allowbreak{}binom(F\textasciicircum{}\allowbreak{}2+\allowbreak{}t,t), without first assuming the factorial identity.



The ceiling rows say F=\allowbreak{}ceil(V0/\allowbreak{}B). Write T=\allowbreak{}t! and



\begin{verbatim}
V0/B=T*product_(i=1)^t (F^2+1)/(F^2+i).
\end{verbatim}

Every factor is \textless{}\allowbreak{}=\allowbreak{}1, so F\textless{}\allowbreak{}=\allowbreak{}T. Monotonicity in F gives V0/\allowbreak{}B\textgreater{}\allowbreak{}=\allowbreak{}2\textasciicircum{}\allowbreak{}t/\allowbreak{}(t+\allowbreak{}1)\textgreater{}\allowbreak{}t\textasciicircum{}\allowbreak{}2 for t\textgreater{}\allowbreak{}=\allowbreak{}12, hence F\textgreater{}\allowbreak{}t\textasciicircum{}\allowbreak{}2. The product bound product(1-a\_\allowbreak{}i)\textgreater{}\allowbreak{}=\allowbreak{}1-sum a\_\allowbreak{}i with a\_\allowbreak{}i=\allowbreak{}(i-1)/\allowbreak{}(F\textasciicircum{}\allowbreak{}2+\allowbreak{}i) gives



\begin{verbatim}
T*(1-1/(2F))<V0/B<=F.
\end{verbatim}

Thus T\textless{}\allowbreak{}F+\allowbreak{}F/\allowbreak{}(2F-1)\textless{}\allowbreak{}=\allowbreak{}F+\allowbreak{}1, so T\textless{}\allowbreak{}=\allowbreak{}F by integrality. Therefore F=\allowbreak{}t!. Conversely at F=\allowbreak{}T the same estimate gives T-1/\allowbreak{}2\textless{}\allowbreak{}V0/\allowbreak{}B\textless{}\allowbreak{}=\allowbreak{}T, furnishing natural ceiling witnesses. This is the inherited square-radius identity.



For 1\textless{}\allowbreak{}=\allowbreak{}i\textless{}\allowbreak{}=\allowbreak{}t, d\_\allowbreak{}i=\allowbreak{}1+\allowbreak{}F\textasciicircum{}\allowbreak{}2/\allowbreak{}i is integral and congruent to one modulo F. Any prime divisor of d\_\allowbreak{}i exceeds t, and the d\_\allowbreak{}i are pairwise coprime, because a common prime for distinct indices would divide their difference. Their product is precisely B. This recovers the reciprocal CRT carrier.



\subsection{Why x is enough amplification}



Retain the literal source majorant with E=\allowbreak{}2\textasciicircum{}\allowbreak{}23 and base R=\allowbreak{}2+\allowbreak{}t+\allowbreak{}x+\allowbreak{}y+\allowbreak{}F, so each nonnegative local source value and local label sum is strictly below R\textasciicircum{}\allowbreak{}E at the prescribed local bounds. The inherited J majorant has ell=\allowbreak{}2\textasciicircum{}\allowbreak{}22.



After F=\allowbreak{}t!, the endpoint equations give



\begin{verbatim}
y>=F+1,
x=1+(1+(t+1)y)*u_t,     x>=F+1>=2.
\end{verbatim}

Thus u\_\allowbreak{}t\textgreater{}\allowbreak{}=\allowbreak{}1 and x\textgreater{}\allowbreak{}=\allowbreak{}2+\allowbreak{}(t+\allowbreak{}1)y. For t\textgreater{}\allowbreak{}=\allowbreak{}12, t\textless{}\allowbreak{}y and F\textless{}\allowbreak{}y, so



\begin{verbatim}
R<2+x+3y<2x,
x>=2+13*(12!+1)>2^32.
\end{verbatim}

At r0=\allowbreak{}2\textasciicircum{}\allowbreak{}32, (2r0)\textasciicircum{}\allowbreak{}E=\allowbreak{}2\textasciicircum{}\allowbreak{}(33E)\textless{}\allowbreak{}2\textasciicircum{}\allowbreak{}r0 because 33E=\allowbreak{}276,824,064\textless{}\allowbreak{}2\textasciicircum{}\allowbreak{}32. For r\textgreater{}\allowbreak{}=\allowbreak{}r0, the binomial theorem and binom(E,j)\textless{}\allowbreak{}=\allowbreak{}E\textasciicircum{}\allowbreak{}j give



\begin{verbatim}
(1+1/r)^E <= sum_(j=0)^E (E/r)^j
             < 1/(1-E/r) < 2,
\end{verbatim}

since E/\allowbreak{}r\textless{}\allowbreak{}=\allowbreak{}1/\allowbreak{}512. Induction therefore proves (2r)\textasciicircum{}\allowbreak{}E\textless{}\allowbreak{}2\textasciicircum{}\allowbreak{}r for all integers r\textgreater{}\allowbreak{}=\allowbreak{}r0. In particular every required local value is



\begin{verbatim}
<R^E<(2x)^E<2^x=2^K.
\end{verbatim}

This proof does not assume a bound on x in terms of y; the arbitrarily large beta representative x makes the estimate stronger, not weaker. For the only tightened ranges, the inherited beta equalities in a canonical finite-refutation witness have u=\allowbreak{}(x-S)/\allowbreak{}(1+\allowbreak{}(i+\allowbreak{}1)y) and v=\allowbreak{}(x-Sp)/\allowbreak{}(1+\allowbreak{}(i+\allowbreak{}2)y), with S,Sp\textgreater{}\allowbreak{}=\allowbreak{}0, y\textgreater{}\allowbreak{}=\allowbreak{}1 and x\textgreater{}\allowbreak{}0. Each denominator is at least two, so u,v\textless{}\allowbreak{}=\allowbreak{}x/\allowbreak{}2\textless{}\allowbreak{}x. Thus completeness admits lower index x instead of x+\allowbreak{}1 for these two ranges. Soundness only strengthens the old bounds; all local values still lie in the old majorant's permitted box.



All lower range indices t,t+\allowbreak{}1,y+\allowbreak{}1,x are \textless{}\allowbreak{}=\allowbreak{}K. Equality for x is permitted: in the completeness valuation estimate division by k! loses at most k-1\textless{}\allowbreak{}=\allowbreak{}K-1, which is still strictly less than K.



\subsection{The self-sized radix and CRT soundness}



Every independently chosen positive representative is \textless{}\allowbreak{}S. For every such representative n and every 0\textless{}\allowbreak{}=\allowbreak{}j\textless{}\allowbreak{}=\allowbreak{}k\textless{}\allowbreak{}=\allowbreak{}K,



\begin{verbatim}
binom(n,j)<=max(1,n^K)<S^K<q.
\end{verbatim}

The coefficients with j\textgreater{}\allowbreak{}n are zero. No aggregate upper-bound equation or slack is required. The literal majorant gives J\textless{}\allowbreak{}=\allowbreak{}S\textasciicircum{}\allowbreak{}ell. At t=\allowbreak{}12, 3ell(t+\allowbreak{}1)\textless{}\allowbreak{}2\textasteriskcentered{}\allowbreak{}12!, and this inequality persists with t by factorial growth. Hence ell(t+\allowbreak{}1)\textless{}\allowbreak{}F\textless{}\allowbreak{}K, and every required coefficient of J through index t+\allowbreak{}1 is also \textless{}\allowbreak{}S\textasciicircum{}\allowbreak{}K\textless{}\allowbreak{}q. These bounds apply to arbitrary representatives in soundness, including the representative deliberately enlarged below.



The modular graphs therefore give all ten intended binomial divisibilities by N. The inherited prime-power lemma says that p\textasciicircum{}\allowbreak{}e|binom(a,k), k,e\textgreater{}\allowbreak{}=\allowbreak{}1, forces a modulo p\textasciicircum{}\allowbreak{}e to lie in [0,k-1]. It is used only in that direction. For row i-1 choose p|d\_\allowbreak{}i and m=\allowbreak{}v\_\allowbreak{}p(d\_\allowbreak{}i). The raw B divisibility of C+\allowbreak{}F\textasciicircum{}\allowbreak{}2+\allowbreak{}1 identifies C modulo p as i-1. Its range then identifies the small residue modulo p\textasciicircum{}\allowbreak{}(Km) as exactly that index. The remaining ranges recover every local coordinate in its prescribed interval. The local source divisibility by N, together with its strict bound \textless{}\allowbreak{}2\textasciicircum{}\allowbreak{}K\textless{}\allowbreak{}=\allowbreak{}p\textasciicircum{}\allowbreak{}(Km), forces the nonnegative local source value to equal zero. The derived-label range similarly forces J(local)\textless{}\allowbreak{}=\allowbreak{}t. The unchanged beta consistency and pairwise-coprime weight/\allowbreak{}Horn-flow argument then produce a finite C222 refutation.



\subsection{Constructive completeness and anchor dominance}



Start with a finite refutation, expand to a proof tree and pad a horizon t\textgreater{}\allowbreak{}=\allowbreak{}12 beyond every required local parameter and label sum. The inherited beta construction supplies x,y and the local source. Set F=\allowbreak{}t!, B, K=\allowbreak{}x and N=\allowbreak{}B\textasciicircum{}\allowbreak{}K. Choose each of the nine positive independent representatives by CRT modulo d\_\allowbreak{}i\textasciicircum{}\allowbreak{}(2K), initially in (0,N\textasciicircum{}\allowbreak{}2]. Replace only A\_\allowbreak{}a by A\_\allowbreak{}a+\allowbreak{}N\textasciicircum{}\allowbreak{}2. This preserves every local polynomial congruence to precision 2K, and leaves every representative \textless{}\allowbreak{}=\allowbreak{}2N\textasciicircum{}\allowbreak{}2.



For each needed range, the falling-factorial numerator contains a term with at least 2Km p-adic valuation, while division by k! loses at most k-1\textless{}\allowbreak{}=\allowbreak{}K-1. The result therefore has at least Km valuation. Polynomial congruence gives the same conclusion for J. This is the genuine precision-2K completeness argument, not the false converse of the range lemma.



Let R=\allowbreak{}2+\allowbreak{}t+\allowbreak{}x+\allowbreak{}y+\allowbreak{}F. The endpoint inequalities above give R\textless{}\allowbreak{}2x. More directly, F,y,t\textless{}\allowbreak{}x, so the remaining nominal indices t,y+\allowbreak{}1,x,F\textasciicircum{}\allowbreak{}2+\allowbreak{}t are \textless{}\allowbreak{}=\allowbreak{}x\textasciicircum{}\allowbreak{}2+\allowbreak{}x. For x\textgreater{}\allowbreak{}=\allowbreak{}5, x\textasciicircum{}\allowbreak{}2+\allowbreak{}x\textless{}\allowbreak{}2\textasciicircum{}\allowbreak{}x, by the base case x=\allowbreak{}5 and induction. Since N=\allowbreak{}B\textasciicircum{}\allowbreak{}x\textgreater{}\allowbreak{}=\allowbreak{}2\textasciicircum{}\allowbreak{}x, every such index is \textless{}\allowbreak{}N. All other nonanchor indices are representatives \textless{}\allowbreak{}=\allowbreak{}2N\textasciicircum{}\allowbreak{}2. Thus s=\allowbreak{}max\_\allowbreak{}i k\_\allowbreak{}i\textless{}\allowbreak{}=\allowbreak{}2N\textasciicircum{}\allowbreak{}2. But



\begin{verbatim}
H=J>=2(A_a+N^2)^2>2N^4>4N^2>=2s,
\end{verbatim}

as N\textgreater{}\allowbreak{}=\allowbreak{}2. This supplies the extra completeness hypothesis of (1), at no additional witness or row cost. Compute S and q afterwards; their carry bounds hold regardless of the actual representatives' sizes. Then compute the exact denominators, W, L and a sufficiently large A satisfying (2). Supply canonical Pell pairs, fresh full-anchor witnesses, Q from (1), and Euclidean modular-digit witnesses. There is no cyclic choice involving q, A, the anchor, or the CRT representatives.



\subsection{Consecutive coordinates and exact contraction}



For every nonanchor seed the chosen endpoint stores natural Y,Z and a product wire P=\allowbreak{}YZ, and imposes



\begin{verbatim}
Y^2+Z^2-2A*P=1,       nY=Z+R_n,
\end{verbatim}

where n is an exact-power base or q+\allowbreak{}1 and R\_\allowbreak{}n is the nonnegative sum in its actual output row. Each seed is checked to have such an incidence. The frame gives 1\textless{}\allowbreak{}=\allowbreak{}n\textless{}\allowbreak{}A. The output row excludes Y=\allowbreak{}0 because the norm would force Z=\allowbreak{}1. It gives Z\textless{}\allowbreak{}=\allowbreak{}nY\textless{}\allowbreak{}AY, so X=\allowbreak{}AY-Z\textgreater{}\allowbreak{}0. Expanding the conic gives X\textasciicircum{}\allowbreak{}2-(A\textasciicircum{}\allowbreak{}2-1)Y\textasciicircum{}\allowbreak{}2=\allowbreak{}1; positive Pell classification follows. Thus (1) applies to its ordinate, and canonical Z=\allowbreak{}Y\_\allowbreak{}(k-1)(A) supplies completeness. The full anchor retains ordinary positive coordinates and every private witness. No anchor norm is silently replaced by an unindexed seed.



Finally the raw polynomial graph is exported and contracted only through explicit integer polynomial restoration certificates. Their nonnegativity is witnessed by rational sums of affine squares plus nonnegative monomials. The independent checker expands these identities, replays each substitution, checks that rows stay quadratic, and independently collects the complete sum of squares. Thus both directions of natural existence are preserved. Finite arithmetic tests and syntax checks complement the ordinary unbounded proof above; they do not replace it or the inherited logical calibration.



\subsection{Fixed-priority 228-variable companion}



The 228-variable companion starts from the identical raw semantic equation. Its contractor uses a fixed, explicitly listed order of row priorities, bound to the raw equation's canonical SHA-256, rather than the predecessor's ascending row order. Every elimination is still an integer monic substitution certified nonnegative by exactly expanded rational squares and nonnegative monomials. The final certificate has 106 eliminations, including two restorations with negative expanded coefficients; both have independently checked positive-square certificates. The target has 228 variables, 178 residuals, 9,866 quartic terms and coefficient height 2,445,412,949,744. No semantic inference or domain relaxation is used for this extra variable reduction. Independent normal and optimized runs replay every restoration and all collected coefficients.
